\section{实验评估}
\label{sec: experiments}

为了进一步理解 TTLM 各变体的性质，我们现在将 TTLM、TTLM-Large 和 TTLM-Tiny 与 Second-order RNNs、RACs、MI-RNNs 以及 Vanilla-RNNs 进行对比，考察它们的有效性。


我们在第 \ref{sec:setting} 节中给出实验设置，在第 \ref{sec:implementations} 节中给出实现细节。我们在第 \ref{sec: rank} 节中研究秩对 TTLM 各变体性能的影响，并在第 \ref{sec: activation} 节中考察非线性激活函数对 TTLM 各变体有效性的影响。



\subsection{实验设置}
\label{sec:setting}
\begin{table*}[t]
\small
\centering
  \begin{tabular}{l|cc|cc|ccc}
    \toprule
      \multirow{2}{*}{模型} & \multicolumn{2}{c|}{
      WikiText-2} &
    \multicolumn{2}{c|}{PTB} & 隐藏单元 & 层数 & 嵌入 \\

    &参数量 & PPL  & 参数量  & PPL & (秩) & & 维度\\
    % \hline

    % RNN \citep{mikolov2012context} & - & - & - & 124.7 & 300 & 1 & - \\
    % LSTM \citep{zaremba2014recurrent} & - & - & - & 114.5 & 200 & 2 & -\\
    % LSTM \citep{grave2016improving} & -& 99.3 & -& 82.3& 1024 & 1 & -\\
    % LSTM \citep{merity2017regularizing} & - & 100.9 & - & 80.6 & 650 & 2 & -\\
    \midrule
    Transformer \citep{vaswani2017attention} & 90.5M & 293.0 & 32.8M & 208.7 &20 & 1 & 400 \\
    Vanilla-RNNs \citep{mikolov2012context} & 11.6M & 96.6 & 4.0M & 115.3 & 20 & 1 & 400\\
    Second-order RNNs \citep{hochreiter1997long} & 11.8M & 96.0 & 4.2M& 108.2 & 20 & 1 & 400\\
    RACs  \citep{levine2018benefits} &11.6M& 97.6 & 4.0M  & 116.8 &20 & 1 & 400 \\
    MI-RNNs \citep{wu_multiplicative_2016} &  11.6M & 99.6 &4.0M & 119.1 & 20 & 1 & 400 \\
    \hline
    TTLM & 12.2M & 546.4 & 4.2M & 559.8 & 20& 1 & 400\\
    TTLM-Tiny & 11.6M &94.9 & 4.0M &106.8 & 20 & 1 & 400\\
    TTLM-Large &11.8M &\textbf{82.3} & 4.2M & \textbf{99.3} & 20 & 1 & 400 \\
  \bottomrule
\end{tabular}
\caption{WikiText-2 与 PTB 数据集上的测试集 PPL。符号 “$-$” 表示原始论文中未提供这些数据。“参数量”列表示参数的数量；详细描述见第 \ref{sec:implementations} 节。“隐藏单元（秩）”列表示隐藏单元或秩的数量。“嵌入维度”列表示每个嵌入向量的维度。我们报告了注意力头数从 [2, 4, 5, 8] 中选取的 Transformer 的最低测试集 PPL。\label{tab:evaluation}}
\end{table*}
%
\paragraph{任务、数据集与指标。} 我们在两个词级语言模型数据集上进行实验：\textbf{(1)} 英文宾州树库（Penn Treebank，PTB）\citep{marcinkiewicz1994building}，由 929k 训练 token、73k 验证 token 和 82k 测试 token 组成，其词表规模为 10k。\textbf{(2)} WikiText-2 数据集 \citep{merity2016pointer} 源自维基百科文章，由 2088k 训练 token、217k 验证 token、45k 测试 token 以及超过 30k 个词类型的词表组成。我们在语言建模任务上比较这些模型，并采用 perplexity（PPL）\citep{meister-cotterell-2021-language} 进行评估；PPL 越低，模型越好。

% \subsubsection{Baselines and Implements}

\paragraph{基线。} 我们的模型与以下基线进行比较：
  Transformer \citep{vaswani2017attention}、Vanilla-RNNs \citep{mikolov2012context}、Second-order RNNs \citep{hochreiter1997long}、循环算术电路（Recurrent Arithmetic Circuits，RACs）\citep{levine2018benefits}，以及乘性积分 RNN（Multiplicative Integration RNNs，MI-RNNs）\citep{wu_multiplicative_2016}。实现细节见第 \ref{sec:implementations} 节。


% \subsubsection{Parameter Setting}
% \label{sec:experimental_setting}

% (2)  Different TT cores settings will influence the number of trainable parameters. Sec. \ref{sec: model size} compares the \textbf{model size} of TT variants

% \begin{comment}

% \subsubsection{Evaluation Metric}
% It  is  the  average  per-word  log-probability on the holdout data set. The low the perplexity, the more effective the model is. The formula of PPL is shown in Eq.\ref{eq:ppl}:

% \begin{equation} \small
% \label{eq:ppl}
%     \text{PPL} = e^{-\frac{1}{n}\sum_{i}\log(p(w_i))}
% \end{equation}
% \end{comment}

% \textbf{Reproductibility statements} To foster reproducibility, all of our code in Sec. \ref{eq: experiments} are provided in the  \textit{link}


% \textbf{Training environments.}

% \textbf{Initialization} The learning weights of embedding tensor $\tW^{xe} \in \mathbb{R}^{|V|\times R\times R}$ were uniformly initialized in the interval $[-0.1,0.1]$ and all other weights were initialized between $[-\frac{1}{\sqrt{H}}, \frac{1}{\sqrt{H}}]$, where $H$ is the hidden units for RNN models. For TTLM-Large or TTLM-Tiny, The learning tensor $\tW^{eh}$ is initialized between $[-\frac{1}{R}, \frac{1}{R}]$ where R is the rank of TTLM. $\mW^{hh}$ is initialized between $[-\frac{1}{\sqrt{R}}, \frac{1}{\sqrt{R}}]$.

\paragraph{超参数。} \textbf{(1)} 为了在相同规模下比较同类模型的有效性，我们将 TTLM 各变体/Vanilla-RNNs 的秩/隐藏单元数设为 [5, 10, 20, 25, 30, 35, 40, 45, 50]。这些模型的嵌入维度为隐藏单元数/秩的平方。这一设置源于第 \ref{sec: Tinyttlm} 节中介绍的 TTLM-Large 与 TTLM-Tiny 的架构。\textbf{(2)} 为避免嵌入维度对 Vanilla-RNNs 性能的潜在影响，我们通过将其嵌入维度设为 [100, 200, 300]，为该模型提供了若干常见的嵌入维度选择。我们将它们相应地命名为 RNNs-100、RNNs-200 和 RNNs-300，并在图 \ref{fig:rank} 中展示。\textbf{(3)} 我们将所有模型训练 50 个 epoch，并选择验证集上的最优模型来预测测试集上的结果。\textbf{(4)} 通过随机梯度下降调整模型中的权重，以最小化训练序列上的平均交叉熵损失，其中梯度采用截断的随时间反向传播（truncated backpropagation through time）算法 \citep{werbos1990backpropagation,williams1990efficient} 计算。随机种子被固定，以确保实验结果不受权重初始化的影响。
